\section{Worum es geht}

\textbf{[SETZUNG]} In der FFGFT werden alle Formeln in natürlichen Einheiten
aufgestellt ($c=\hbar=k_B=1$). Beim Rückrechnen auf SI-Größen erscheinen die
physikalischen Konstanten mit bestimmten Potenzen wieder. Diese Potenzen sind
\emph{nicht} frei wählbar -- sie folgen zwingend aus der Dimensionsanalyse.
Ein Fehler bei der Potenz von $c$ ist ein Fehler von Faktoren
$\sim3\times10^8$ bis $(3\times10^8)^4$ -- nicht bemerkbar wenn man nur auf
Zehnerpotenzen schaut, fatal wenn man das Ergebnis mit einer Messung vergleicht.

Dieses Dokument gibt das systematische Verfahren: erst Dimensionsanalyse in
natürlichen Einheiten, dann maschinelle Rückrechnung, dann Probe. Es ist kein
Lehrdokument über Einheitensysteme (das ist A080/A085) -- es ist ein
Handwerks-Werkzeugkasten für den konkreten Rechenfall.

\section{Grundprinzip: alle Dimensionen als Energiepotenzen}

\textbf{[B]} In natürlichen Einheiten gilt $c=\hbar=k_B=1$. Damit hat jede
physikalische Größe eine eindeutige Dimension $[E]^n$ mit rationalem $n$:

\begin{center}
\renewcommand{\arraystretch}{1.35}
{\small
\begin{tabular}{llll}
\toprule
Größe & Nat.~Dim. & SI-Herkunft & $n$ \\
\midrule
Energie & $[E]^1$ & --- & $+1$ \\
Masse & $[E]^1$ & $E=mc^2\Rightarrow m=E/c^2$ & $+1$ \\
Impuls & $[E]^1$ & $p=E/c$ & $+1$ \\
Länge & $[E]^{-1}$ & $\ell=\hbar c/E$ & $-1$ \\
Zeit & $[E]^{-1}$ & $t=\hbar/E$ & $-1$ \\
Geschwindigkeit & $[E]^0$ & $v=\ell/t=c=1$ & $0$ \\
Kraft & $[E]^2$ & $F=E/\ell=E^2/(\hbar c)$ & $+2$ \\
Druck/Energiedichte & $[E]^4$ & $\rho=E/\ell^3=E^4/(\hbar c)^3$ & $+4$ \\
Ladung & $[E]^0$ & aus $\alpha=e^2/(4\pi\varepsilon_0\hbar c)$ dimensionslos & $0$ \\
Temperatur & $[E]^1$ & $T=E/k_B$ & $+1$ \\
Entropie & $[E]^0$ & $S=k_B\cdot(\text{dimensionslos})$ & $0$ \\
Gravitationskonstante $G$ & $[E]^{-2}$ & $G=\ell_P^2c^3/\hbar$ & $-2$ \\
\bottomrule
\end{tabular}
}
\renewcommand{\arraystretch}{1.0}
\end{center}

\textbf{[B]} \textbf{Merksatz.} Aus $n$ folgt der Rückrechnungsfaktor direkt:
eine Größe mit $[E]^n$ in natürlichen Einheiten hat den SI-Wert
\[
  X_\text{SI} = X_\text{nat}\cdot\left(\frac{\hbar c}{1\,\text{SI-Länge}}\right)^n.
\]
In der Praxis: man stellt den SI-Faktor aus $\hbar$, $c$, $k_B$ so zusammen,
dass die SI-Dimension herauskommt.

\section{Die c-Potenz beim Rückrechnen: systematische Herleitung}

\textbf{[B]} Die Frage „welche Potenz von $c$?" beantwortet sich aus der
SI-Dimension der Zielgröße. Die Methode in drei Schritten:

\textbf{Schritt 1: Zieldimension in SI aufschreiben.}
z.B.\ Energiedichte: $[\rho]=\text{J/m}^3=\text{kg\,m}^{-1}\text{s}^{-2}$.

\textbf{Schritt 2: SI-Dimension durch $\hbar^a c^b k_B^d$ ausdrücken.}
$\hbar$ hat SI-Dim $[\text{J\,s}]$, $c$ hat $[\text{m/s}]$, $k_B$ hat $[\text{J/K}]$.
Gesucht: $a,b,d\in\mathbb Z$ so dass $[\hbar^a c^b k_B^d]=$ Zieldim/$[E_\text{nat}^n]$.

\textbf{Schritt 3: Probe durch Einsetzen.}
Mit dem gefundenen $(a,b,d)$ einsetzen und auf die gemessene SI-Größe prüfen.

\textbf{[K]} \textbf{Tabelle der häufigsten Rückrechnungsfaktoren:}

\begin{center}
\renewcommand{\arraystretch}{1.35}
{\small\setlength\tabcolsep{4pt}
\begin{tabular}{lllp{4.5cm}}
\toprule
Größe & nat.\ Einheit & SI-Faktor & Rechenweg \\
\midrule
Masse & $[E]$ in MeV & $\div c^2$ & $m_e=0{,}511\,\text{MeV}/c^2$ \\
Länge & $[E]^{-1}$ & $\times\hbar c$ & $\ell_P=1/E_P\cdot\hbar c$ \\
Zeit & $[E]^{-1}$ & $\times\hbar$ & $t_P=1/E_P\cdot\hbar$ \\
Geschwindigkeit & dimensionslos & $\times c$ & $v=0{,}5\cdot c$ \\
Kraft & $[E]^2$ & $\div(\hbar c)$ & $F[\text{N}]=F_\text{nat}\cdot E^2/(\hbar c)$ \\
Energiedichte & $[E]^4$ & $\div(\hbar c)^3$ & $\rho[\text{J/m}^3]=\rho_\text{nat}\cdot E^4/(\hbar c)^3$ \\
Temperatur & $[E]$ & $\div k_B$ & $T[\text{K}]=E[\text{J}]/k_B$ \\
Frequenz $f$ & $[E]^1$ & $\div\hbar$ & $f[\text{s}^{-1}]=E[\text{J}]/\hbar$ \\
$G$ & $[E]^{-2}$ & $\times\hbar c^3$ & $G[\text{m}^3/\text{kg\,s}^2]=G_\text{nat}\cdot\hbar c^3$ \\
\bottomrule
\end{tabular}
}
\renewcommand{\arraystretch}{1.0}
\end{center}

\section{Konkrete Beispiele aus der FFGFT}

\subsection{Masse: eine Potenz von $c^2$}

\textbf{[K]} Aus $E=mc^2$ folgt $m=E/c^2$. Ein Wert in natürlichen Einheiten
(MeV) wird zur SI-Masse durch Division durch $c^2$:
\[
  m_e = 0{,}511\,\frac{\text{MeV}}{c^2}
  = \frac{0{,}511\times10^6\times1{,}602\times10^{-19}\,\text{J}}{(2{,}998\times10^8\,\text{m/s})^2}
  = 9{,}109\times10^{-31}\,\text{kg}.
\]
Fehler $c^1$ statt $c^2$: Faktor $3\times10^8$ daneben.
Fehler $c^0$ (kein $c$): Faktor $(3\times10^8)^2\approx10^{17}$ daneben.

\subsection{Planck-Länge: $\hbar c$ zusammen, nicht $c$ allein}

\textbf{[K]} $\ell_P$ in natürlichen Einheiten ist eine
inverse Energie. Die Rückrechnung auf Meter braucht $\hbar c$:
\[
  \ell_P = \frac{\hbar c}{E_P}
  = \frac{(1{,}055\times10^{-34}\,\text{J\,s})(2{,}998\times10^8\,\text{m/s})}
         {1{,}956\times10^9\,\text{J}}
  = 1{,}616\times10^{-35}\,\text{m}.
\]
Nur $\hbar/E_P$ (ohne $c$) ergibt die Planck-\emph{Zeit}, nicht die
Planck-\emph{Länge} -- ein typischer Fehler.

\subsection{Energie zu inverser Zeit: nur $\hbar$, kein $c$}

\textbf{[K]} Eine Energie $E$ (Dimension $[E]^1$) wird zur inversen Zeit
$[\text{s}^{-1}]$ durch Division durch $\hbar$ zurückgerechnet:
\[
  f = \frac{E}{\hbar},
  \qquad
  [f] = \frac{[\text{J}]}{[\text{J\,s}]} = [\text{s}^{-1}].\ \checkmark
\]
Beispiel: $E=1{,}41\times10^{-33}\,\text{eV}=2{,}26\times10^{-52}\,\text{J}$
ergibt $f=E/\hbar=2{,}14\times10^{-18}\,\text{s}^{-1}$ (A265).
Würde man irrtümlich durch $\hbar c$ dividieren, käme eine inverse Länge
(falsche Dimension). Die $c$-Potenz ist null -- sie erscheint nicht.

\subsection{Energiedichte: vier Potenzen von $c$ im Nenner}

\textbf{[K]} Casimir-Energiedichte $\rho$ hat Dimension $[E]^4$ in natürlichen
Einheiten. Rückrechnung:
\[
  \rho[\text{J/m}^3]
  = \rho_\text{nat}\,[E^4]\cdot\frac{1}{(\hbar c)^3}.
\]
Drei Potenzen von $\hbar c$ im Nenner -- weil Länge dreimal $[E]^{-1}$ liefert.
$c$ erscheint als $c^3$ (aus den drei $\hbar c$).

\subsection{Gravitationskonstante: $c^3$ im Zähler}

\textbf{[K]} In natürlichen Einheiten hat $G$ Dimension $[E]^{-2}$. Die
SI-Konstante $G=6{,}674\times10^{-11}\,\text{m}^3\text{kg}^{-1}\text{s}^{-2}$
erfordert:
\[
  G_\text{SI} = G_\text{nat}\cdot\hbar c^3,
  \qquad
  [G_\text{SI}] = [E]^{-2}\cdot[\hbar][c]^3
  = \frac{\text{J\,s}\cdot\text{m}^3/\text{s}^3}{\text{J}^2}
  = \frac{\text{m}^3}{\text{kg\,s}^2}.\ \checkmark
\]
Fehler: $c^2$ statt $c^3$ ergibt falsche Dimension sofort erkennbar.

\section{Das allgemeine Verfahren: Dimensionsmatrix}

\textbf{[B]} Für jede Formel $X=f(\xi,E_0,\ldots)$ in natürlichen Einheiten:

1. Bestimme $n$ aus der Tabelle oben (oder durch Zählen der
   Längen/Massen/Zeiten in der SI-Dimension von $X$).

2. Schreibe den SI-Faktor:
\[
  X_\text{SI} = X_\text{nat}\cdot\hbar^a\cdot c^b\cdot k_B^d
\]
mit $a,b,d$ aus dem linearen System:
\begin{align*}
  [\text{kg}]:&\quad -a+b = n_\text{kg} \\
  [\text{m}]:&\quad a+b = n_\text{m} \\
  [\text{s}]:&\quad -a-b = n_\text{s}
\end{align*}
(wobei $n_\text{kg}$, $n_\text{m}$, $n_\text{s}$ die Exponenten der SI-Dimension
von $X$ sind; $k_B$ kommt hinzu wenn Kelvin erscheint.)

3. Probe: setze $\hbar=1{,}055\times10^{-34}\,\text{J\,s}$,
   $c=2{,}998\times10^8\,\text{m/s}$,
   $k_B=1{,}381\times10^{-23}\,\text{J/K}$ ein und prüfe die SI-Einheit.

\textbf{[K]} \textbf{Wann $k_B$ erscheint und wann nicht.} $k_B$ erscheint genau
dann, wenn die Zielgröße Kelvin enthält (Temperatur, Entropie/Kelvin) oder wenn
man von einer Energie auf eine Kelvin-Angabe umrechnet. In allen anderen Fällen
(Massen, Längen, Zeiten, Kräfte) erscheint $k_B$ nicht.

\section{Einheitenprüfung einer Formel}

\textbf{[B]} Jede in FFGFT aufgestellte Formel sollte die Einheitenprüfung
bestehen: beide Seiten der Gleichung müssen in natürlichen Einheiten dieselbe
Energiepotenz $[E]^n$ haben, und beim Rückrechnen muss dieselbe SI-Einheit
entstehen.

\textbf{[K]} \textbf{Beispiel: $T\cdot E = 1$} (A060). In natürlichen Einheiten:
$[T]=[E]^{-1}$ (Zeit) und $[E]=[E]^{+1}$. Produkt: $[E]^0$ -- dimensionslos.
Das Produkt $T\cdot E=1$ ist eine korrekte Gleichung zwischen einer Zeit und
einer inversen Zeit -- die Einheitenprüfung besteht.

\textbf{[K]} \textbf{Beispiel: $\alpha=\xi\cdot E_0^2$ -- Einheitenprüfung der
$\alpha$-Formel.} Die Formel ist \emph{richtig}. Die Einheitenprüfung zeigt
warum -- und macht gleichzeitig die implizite Referenzskala sichtbar.

$\alpha$ ist dimensionslos, $\xi$ ist dimensionslos. Damit stimmt die Gleichung
nur wenn auch $E_0^2$ dimensionslos ist. Das ist der Fall wenn $E_0$ als reines
Zahlenverhältnis erscheint -- also in MeV gemessen und durch $1\,\text{MeV}$
geteilt:
\[
  \alpha = \xi\cdot\Bigl(\frac{E_0}{1\,\text{MeV}}\Bigr)^2.
\]
Die $1\,\text{MeV}$ im Nenner ist die \emph{implizite Referenzskala} -- sie
macht $E_0/1\,\text{MeV}$ dimensionslos und die Gleichung korrekt. Sie steht
nicht explizit in der Kurzschreibweise, ist aber unverzichtbar.

\textbf{Numerische Probe:}
\[
  \alpha = \frac{4}{30000}\cdot\Bigl(\frac{7{,}398\,\text{MeV}}{1\,\text{MeV}}\Bigr)^2
  = 1{,}333\times10^{-4}\times54{,}73
  = 7{,}297\times10^{-3}
  = \frac{1}{137{,}0}.
\]
Übereinstimmung mit dem Experiment auf $0{,}0005\,\%$ -- die Formel stimmt.

\textbf{Was die Einheitenprüfung aufdeckt:} Die Wahl MeV als Referenzskala ist
nicht frei -- sie ist der empirische SI-Anker des FFGFT-Sektors. Eine andere
Referenzskala (z.B.\ GeV) würde $\xi$ einen anderen numerischen Wert geben,
aber $\alpha$ bleibt dieselbe physikalische Zahl. Das ist die Aussage aus A261:
$E_0=7{,}398\,\text{MeV}$ trägt den Anker, und deshalb darf man in dieser
Formel nicht einfach andere Energieeinheiten einsetzen ohne $\xi$ mitzuändern.

\textbf{[S]} \textbf{Umstellung auf natürliche Einheiten.} Eine SI-Formel in
natürliche Einheiten überführen: alle $\hbar$, $c$, $k_B$ durch~$1$ ersetzen
und die Einheiten weglassen. Die resultierende Formel gilt dann nur in natürlichen
Einheiten; beim Rückrechnen kommen alle Faktoren wieder.

\section{Zeichenerklärung}

Symbole die in der gesamten A-Serie verwendet werden, sind in A010 erklärt.
Spezifisch für dieses Dokument:

\begin{center}
\renewcommand{\arraystretch}{1.3}
{\small\begin{tabular}{lp{9cm}}
\toprule
Symbol & Bedeutung \\
\midrule
$[X]$ & Dimension der Größe $X$ \\
$[E]^n$ & Energiedimension zur Potenz $n$ \\
$n_\text{kg},n_\text{m},n_\text{s}$ & Exponenten der SI-Einheit von $X$
  in $\text{kg}^{n_\text{kg}}\text{m}^{n_\text{m}}\text{s}^{n_\text{s}}$ \\
$a,b,d$ & Exponenten von $\hbar^a c^b k_B^d$ im Rückrechnungsfaktor \\
$E_P = \sqrt{\hbar c^5/G}\approx1{,}956\times10^9\,\text{J}$ & Planck-Energie \\
\bottomrule
\end{tabular}}
\renewcommand{\arraystretch}{1.0}
\end{center}

\section{Prüfskript}

\begin{itemize}[leftmargin=2em,itemsep=2pt,topsep=2pt]
  \item Numerische Verifikation aller Rückrechnungsbeispiele;
    Dimensionsmatrix-Löser für beliebige SI-Zieldimensionen;
    Probe für $m_e$, $\ell_P$, inverse Zeit, $\rho_\text{Casimir}$, $G$:
    \texttt{python/A\_Serie\_Skripte/a266\_einheitenpruefung.py}
\end{itemize}

\section{Was offen bleibt}

\textbf{Die Dimensionsmatrix setzt voraus, dass die Formel in reinen
Planck-Einheiten aufgestellt ist.} Mischsysteme (z.B.\ MeV für Energie, fm
für Länge) erfordern zusätzliche Umrechnungsschritte.

\textbf{Nicht-ganzzahlige Energiepotenzen.} Einige FFGFT-Ausdrücke haben
rationale $n$ (z.B.\ Ladung $[E]^{1/2}$). Das Verfahren gilt prinzipiell,
aber das lineare System liefert dann rationale Lösungen.

\section{Ersetzt}

Dieses Dokument nimmt auf: Dok.~013 (Dimensionsanalyse, Gravitationskonstante
aus $\xi$, Planck-Länge, Umrechnungsfaktoren), Dok.~014 (natürliche Einheiten,
Dimensionstabelle, Skalierungsfaktor $S_{T0}$, Herleitung vs.\ Rückrechnung),
Dok.~015 (Einheitenumwandlungen, Energiedimension-Tabelle, Kraft/Druck/Ladung).
\emph{Nicht} übernommen: pädagogische Herleitung der einzelnen Konstanten aus
$\xi$ (A015, A145) und die Einheitenkonvention als solche (A080, A085).

\section{Verwendung von KI-Werkzeugen}

Dieses Dokument ist unter Verwendung KI-gestützter Werkzeuge entstanden.
Verantwortung für Inhalt und Fehler liegt beim Autor. Wortlaut der Regeln in A010.
